\documentclass[aps,prl,reprint,amsmath,amssymb,nofootinbib,raggedbottom]{revtex4-2}
\usepackage[T1]{fontenc}
\usepackage{tnnotes}
\usepackage[colorlinks=true,citecolor=blue,linkcolor=blue,urlcolor=blue]{hyperref}
\hypersetup{pdftitle={Scope of the Spacetime Bootstrap}}

\newtheorem{corollary}{Corollary}

% ------------------------------------------------ diagrams specific to this note
% One bond event: C on the left site, Theta(C) on its mirror image; the dashed
% line is the mirror, the double lines are the two sites' imaginary-time legs.
\newcommand{\bondform}[2]{\tikz[tn,tn centered]{
 \draw[phys] (0,-1)--(0,1); \draw[phys] (3.2,-1)--(3.2,1);
 \draw (0,0)--(3.2,0);
 \draw[black!45,densely dashed,line width=.35pt] (1.6,-1.15)--(1.6,1.15);
 \node[xop,fill=white] at (0,0) {$#1$};
 \node[xop,fill=white] at (3.2,0) {$#2$};}}
% A transfer vertex with named virtual legs.
\newcommand{\Xab}[2]{\tikz[tn]{\draw (-.95,0)--(.95,0); \node[xvertex] at (0,0) {};
 \node[anchor=east] at (-1,0) {$#1$}; \node[anchor=west] at (1,0) {$#2$};}}
% One site of the transfer vertex, opened in imaginary time: the physical
% trace (closed double loop) with the left-bond insertions #1 and the
% right-bond insertions #2.
\newcommand{\eventsite}[2]{\tikz[tn,tn centered]{
 \begin{pgfonlayer}{background}
 \draw[phys,ring] (0,-1.9)--(0,1.9)--(1.9,1.9)--(1.9,-1.9)--cycle;
 \draw (-1.9,1)--(0,1); \draw (-1.9,-1)--(0,-1);
 \draw[preaction={draw,white,line width=2.5pt}] (0,0)--(2.6,0);
 \end{pgfonlayer}
 \node[xop,fill=white] at (0,1) {$#1$};
 \node[xop,fill=white] at (0,-1) {$#1$};
 \node[xop,fill=white] at (0,0) {$#2$};
 \node[anchor=east] at (-1.9,0) {$a$}; \node[anchor=west] at (2.6,0) {$b$};}}
% The reflection-positivity inequality: a spacetime torus cut by the mirror,
% A on the left half, theta(A) on its image.
\newcommand{\rpcut}{\tikz[tn,tn centered]{
 \fill[mpoblue!12] (0,0) rectangle (4.4,1.8);
 \draw[black!60,line width=.35pt] (0,0) rectangle (4.4,1.8);
 \draw[black!45,densely dashed,line width=.35pt] (2.2,-.25)--(2.2,2.05);
 \node[xop,fill=white] at (1.1,.9) {$A$};
 \node[xop,fill=white] at (3.3,.9) {$\theta(A)$};}}
% delta(eps) = [(1-eps)^-3 - 1]^2/2 against the diagonal, 0 <= eps <= .16;
% they cross at eps* = .127.
\newcommand{\deltacurve}{\tikz[baseline=0pt,x=6.5em,y=6.5em,line width=.45pt]{
 \draw[black!35,line width=.3pt] (0,0)--(.16,0) (0,0)--(0,.16);
 \draw[black!45,densely dotted] (0,0)--(.16,.16);
 \draw[mpored,domain=0:.1405,samples=30,smooth]
   plot (\x,{.5*(1/((1-\x)*(1-\x)*(1-\x))-1)*(1/((1-\x)*(1-\x)*(1-\x))-1)});
 \fill (.1273,.1273) circle[radius=.035em];}}
% A staggered longitudinal field along a chain, matching \fieldsame and \fieldopp.
\newcommand{\fieldstag}{\tikz[baseline=0pt,x=1em,y=1em,line width=.5pt]{
 \draw[black!55] (0,.4)--(1.8,.4);
 \foreach \x/\s in {0/1,.6/-1,1.2/1,1.8/-1}{
   \fill (\x,.4) circle[radius=.08];
   \draw[-{Latex[length=.95mm,width=.8mm]},mpored,line width=.6pt]
     (\x,{.4-.36*\s})--(\x,{.4+.36*\s});}}}
% Two-site unit cell with internal swap F across an inter-cell mirror cut.
\newcommand{\cellbondform}{\tikz[tn,tn centered]{
 \draw[rounded corners=1.5pt,fill=mpoblue!10,draw=black!45] (-.95,-.65) rectangle (1.15,.65);
 \draw[rounded corners=1.5pt,fill=mpoblue!10,draw=black!45] (2.45,-.65) rectangle (4.55,.65);
 \draw[phys] (-.55,-1)--(-.55,1); \draw[phys] (.65,-1)--(.65,1);
 \draw[phys] (2.95,-1)--(2.95,1); \draw[phys] (4.15,-1)--(4.15,1);
 \draw (.65,0)--(2.95,0);
 \draw[black!45,densely dashed,line width=.35pt] (1.8,-1.15)--(1.8,1.15);
 \node[xop,fill=white] at (.65,0) {$C_\alpha$};
 \node[xop,fill=white] at (2.95,0) {$F\overline{C_\alpha}F$};
 \node[font=\tiny,black!65] at (-.55,0) {$a$};
 \node[font=\tiny,black!65] at (4.15,0) {$b$};}}

\begin{document}
\title{Scope of the Spacetime Bootstrap}
\date{October 15, 2026}
\begin{abstract}
The spacetime bootstrap establishes a uniform spectral gap from bond reflection positivity and certified finite-scale purity bounds. Bond reflection positivity ensures that the continuous-time spatial transfer operator is self-adjoint, enabling spatial scale doubling dual to imaginary time. We formalize the method as a general criterion for one-dimensional translation-invariant lattice systems, extend it to multi-site unit cells with internal reflection and odd seed lengths, classify standard spin and ladder Hamiltonians, and summarize the computer-verified gap certificates across eight model families.
\end{abstract}
\maketitle
\setlength{\abovedisplayskip}{5pt plus 2pt minus 1pt}
\setlength{\belowdisplayskip}{5pt plus 2pt minus 1pt}

\appendix
\setcounter{secnumdepth}{2}
\setcounter{section}{3}

\section{Prerequisites of the Spacetime Bootstrap}\label{app:uses}
The scale-doubling iteration in Ref.~\cite{main} comprises three elementary operations:
\begin{itemize}
\item \emph{Coupling identities.} Equalities between partition-function ratios on the spacetime torus, such as $\Zt42/\Ztsq[4]21$. These hold universally for every translation-invariant lattice Hamiltonian.
\item \emph{Temporal squaring.} The relation $\tncolumn{\rho_n,\rho_n}\propto\rho_n(2\beta)$ requires only positivity of the Gibbs operator, $e^{-\beta H}=e^{-\beta H/2}(e^{-\beta H/2})^\dagger\ge0$, which holds generally by Euclidean time reflection.
\item \emph{Spatial squaring.} The dual relation $\tnchain{\sigma_n,\sigma_n}\propto\sigma_{2n}$ requires positivity of the spatial transfer operator, $\sigma_n=X_\beta^n/Z_n\ge0$. This holds whenever $X_\beta$ is self-adjoint ($X_\beta=X_\beta^\dagger$), which additionally guarantees that $Z_n=\sum_i\lambda_i^n$ with real eigenvalues $\lambda_i$, permitting $\ell^p$ norm interpolation to non-dyadic intermediate lengths.
\end{itemize}
Self-adjointness of $X_\beta$ represents the spatial counterpart of Euclidean time reflection, and is not automatic. Once it holds, the remaining doubling argument follows directly.
The three hypotheses are therefore:
\begin{itemize}
\item[(H1)] \emph{Locality.} A translation-invariant chain with finite local
dimension and nearest-neighbor terms. Finite-range interactions reduce to
this by blocking sites into unit cells.
\item[(H2)] \emph{Bond reflection positivity} (Sec.~\ref{app:rp}).
\item[(H3)] \emph{Input.} Both purities are near one at a single scale
(Sec.~\ref{app:criterion}).
\end{itemize}
Hypotheses (H1) and (H2) are algebraic properties of the Hamiltonian, whereas (H3) reflects the underlying spectral gap.

\section{Reflection positivity}\label{app:rp}
Let $\Theta$ be an antiunitary map on one site with $\Theta^2=\pm1$, for
example complex conjugation $K$ in a chosen basis. For an operator $C$ on one site write
$\Theta(C)=\Theta C\Theta^{-1}$.
\begin{itemize}
\item \emph{Reflection.} Cut the chain between two neighboring sites. Mirror
the chain across the cut and apply $\Theta$ on every site. Call the result
$\theta$. It is antilinear, it maps the operators on the left half to those
on the right half, and $\theta^2=1$.
\item \emph{Bond form.} A nearest-neighbor Hamiltonian is in \emph{bond
reflection-positive form}~\cite{dls,fils,biskup,jaffejanssens} if it reads
\begin{equation}
 H=\sum_jD_j-\sum_{j,k}c_k\,C_{k,j}\,\Theta(C_k)_{j+1},\qquad c_k>0,
 \label{eq:rpform}
\end{equation}
with $\Theta(D)=D$. Each bond event of weight $c_k$ inserts $C_k$ on one
site and its mirror image $\Theta(C_k)$ on the next site,
\begin{equation*}
 -h_{j,j+1}=\sum_kc_k\;\bondform{C_k}{\Theta(C_k)}.
\end{equation*}
The sign in front matters. A term with $c_k<0$ cannot be absorbed by
rescaling $C_k$, because $\Theta$ conjugates the factor $i$:
$(iC)\otimes\Theta(iC)=C\otimes\Theta(C)$.
\end{itemize}
The model of Ref.~\cite{main} has this form with $\Theta=K$ in the Cartesian
basis and $c_k=1$. Its generators $G_\alpha=iS^\alpha$ are real, so
$-\bm S\cdot\bm S=\sum_\alpha G_\alpha\otimes\overline{G_\alpha}$.

\emph{Reflection positivity and transfer matrices.} Expand $e^{-\beta H}$ in the bond
events that cross the mirror. Each term is a left half followed by its mirror
image, so for every operator $A$ on the left half,
\begin{equation}
 \Tr\bigl[A\,\theta(A)\,e^{-\beta H}\bigr]=\rpcut
 =\sum\langle\Phi_A|\Phi_A\rangle\ge0 .
 \label{eq:rp}
\end{equation}
Reflection positivity began as the Osterwalder--Schrader axiom that
rebuilds a Hilbert space and a positive transfer matrix from Euclidean
data~\cite{os}. On the lattice it gives a self-adjoint, positive transfer
matrix~\cite{luscher}, and its Cauchy--Schwarz structure drives infrared
bounds and chessboard estimates~\cite{fss,fils,dls}. For spins and
Majoranas, Jaffe and Janssens characterize exactly which couplings are
reflection positive~\cite{jaffejanssens}. The spacetime bootstrap needs only
the weaker consequence that follows: the network cut at a single bond is
Hermitian. This is the L\"uscher statement with space and time exchanged.

\begin{lemma}[Mirror is conjugation]\label{lem:mirror}
Under Eq.~\eqref{eq:rpform}, build the continuous-time MPO as in
Ref.~\cite{main}. The events are $(\tau,k)$, they occur at rate $c_k$, and
$e^{-\tau D}$ propagates between them. Then the transfer kernel obeys
\begin{equation}
 \overline{\Xab ab}=\Xab ba,\qquad\text{i.e.}\quad
 \overline{X_\beta(a,b)}=X_\beta(b,a).
 \label{eq:mirror}
\end{equation}
\end{lemma}
\begin{proof}
Open the vertex in imaginary time. Events on the left bond $a$ insert
$\Theta(C)$, and events on the right bond $b$ insert $C$, along a single traced
physical loop:
\begin{equation*}
 \overline{X_\beta(a,b)}=\overline{\eventsite{\Theta C}{C}}
 =\eventsite{C}{\Theta C}=X_\beta(b,a).
\end{equation*}
Write $\Theta=VK$. Complex conjugation of a trace conjugates every
insertion and keeps the time order. Conjugating the loop by $V$ then maps
$\overline{\Theta(C)}\mapsto C$, $\overline C\mapsto\Theta(C)$ and
$\overline D\mapsto D$, because $V\overline V=\pm1$ and $\Theta(D)=D$. The
picture that results is the mirror image of the original: the left
insertions are now $C$ and the right ones are $\Theta(C)$.
\end{proof}

\begin{theorem}[Space behaves like time]\label{thm:space}
Assume (H1) and (H2). Then $X_\beta$ is a self-adjoint Hilbert--Schmidt
operator on $L^2(\mu_\beta)$. For every even $n$, $Z_n(\beta)=\sum_i\lambda_i(\beta)^n$ with
real $\lambda_i$, and $\sigma_n(\beta)\ge0$. In particular the squaring
lemma, the coupling identities and the $\ell^p$ interpolation of
Ref.~\cite{main} hold verbatim.
\end{theorem}
\begin{proof}
The event measure $\mu_\beta$ is positive because $c_k>0$, and it is finite,
with mass $e^{\beta\sum_kc_k}$ per bond. Normalize $\|C_k\|=1$. The kernel is
then bounded, $|X_\beta(a,b)|\le d\,e^{\beta\|D\|}$, hence Hilbert--Schmidt,
and it is Hermitian by Lemma~\ref{lem:mirror}. The remaining statements are
those of Ref.~\cite{main} for a self-adjoint $X_\beta$.
\end{proof}

\emph{Multi-site unit cells and internal reflection.}
When $c$ sites are blocked into a unit cell on a ring of $L$ cells,
$H_L=\sum_{j=1}^L h_{j,j+1}+\sum_{j=1}^L o_j$, a mirror cut across an
inter-cell bond swaps the left and right virtual legs of each cell while
reversing the site order inside the cell. Let $F$ be a real symmetric
involution ($F=F^{\mathsf T}$, $F^2=I$) that reverses the internal site order.
If the inter-cell bond $h$ and on-site cell operator $o$ satisfy
\begin{equation}
 -h=\sum_\alpha C_\alpha\otimes F\overline{C_\alpha}F,\qquad
 FoF=o,\qquad o=\overline{o},
 \label{eq:rpcell}
\end{equation}
then $\Theta(C)=F\overline{C}F$ is antiunitary with $\Theta^2=1$ and
$\Theta(o)=o$. Each inter-cell bond event pairs site $b$ of the left cell
with its reflected partner $a$ in the right cell,
\begin{equation*}
 -h_{j,j+1}=\sum_\alpha\;\cellbondform\,,
\end{equation*}
and conjugating the physical loop by $F$ in Lemma~\ref{lem:mirror} gives
$\overline{X_\beta(a,b)}=X_\beta(b,a)$ on the ring of $L$ cells (\texttt{RPRing}
in \path{lean/CertifiedGaps/Ring/Transfer.lean} and \path{Transfer/}). For a
two-spin dimer $(a,b)$ with swap $F=P$, $C_\alpha=I\otimes G_\alpha$ gives
$F\overline{C_\alpha}F=\overline{G_\alpha}\otimes I$, linking spin~$b$ of
cell~$j$ to spin~$a$ of cell~$j+1$.

\emph{Positive cell-transfer operators and odd seed lengths.}
Theorem~\ref{thm:space} restricts $Z_n(\beta)=\sum_i\lambda_i(\beta)^n$ to
even $n$ because a self-adjoint $X_\beta$ may have negative eigenvalues. In a
two-spin dimer cell $(a,b)$ whose intra-cell bond is also reflection positive
with real factors, cutting the cell across its internal bond factors the
cell transfer operator into a half-cell product,
\begin{equation}
 X_\beta=\tnchain[node=xop]{K,K^*}=KK^*\succeq0 .
 \label{eq:cellpos}
\end{equation}
Every transfer eigenvalue $\lambda_i(\beta)$ is then nonnegative, so
$\sigma_n(\beta)=X_\beta^n/Z_n(\beta)\ge0$, the squaring map, and $\ell^p$
interpolation hold at odd as well as even cell counts $n_0$. In the
dimerized spin-1 chain (\path{models/spin1_dimerized}), Eq.~\eqref{eq:cellpos}
enables an odd seed of $n_0=3$ dimer cells ($6$ and $12$ spins), bypassing the
$2n_0=8$-cell ($16$-spin, $3^{16}=4.3\times10^7$-state) ring required by
$n_0=4$.

\emph{Relation to the sign problem.} A Hamiltonian is stoquastic if, in some basis, every
configuration of a quantum Monte Carlo expansion has a nonnegative weight.
Reflection positivity asks less of individual configurations and more of their
structure: a weight may be negative, provided each configuration paired with
its mirror image contributes $|\cdot|^2$.
\begin{itemize}
\item In Ref.~\cite{main} the event measure is positive, but the entries of
$W$ and $X_\beta$ carry signs, because the $G_\alpha$ are real antisymmetric
matrices. A sampler over histories would meet a sign problem; the spectral
argument does not, since $\sum_i\lambda_i^n>0$.
\item Neither property implies the other: the Heisenberg ferromagnet is
stoquastic but not reflection positive~\cite{speer} (Proposition~\ref{prop:xyz}).
\item Sign-free fermion Monte Carlo methods based on time reversal or Majorana
symmetry use the same mirror-squared factorization
$\det M\cdot\overline{\det M}$~\cite{wuzhang,wei,li,jaffe}.
\end{itemize}

\section{The general criterion}\label{app:criterion}
Define the one-step loss
\begin{equation}
 \delta(\varepsilon)=\tfrac12\bigl[(1-\varepsilon)^{-3}-1\bigr]^2 .
 \label{eq:delta}
\end{equation}
It lies below the diagonal, \deltacurve, up to $\varepsilon_*\simeq0.1273$,
the root of $\delta(\varepsilon)=\varepsilon$. Ref.~\cite{main} uses the
cruder bound $\delta\le5\varepsilon^2$, valid for $\varepsilon\le1/100$.

\begin{theorem}[Spacetime-bootstrap criterion]\label{thm:criterion}
Assume (H1) and (H2), and suppose that for some even $n_0$ and some $\beta_0>0$,
\begin{equation}
 \Ps(n_0,\beta_0),\ \Pt(2n_0,\beta_0)\ge1-\varepsilon_0,\qquad
 \varepsilon_0<\varepsilon_* .
 \label{eq:H3}
\end{equation}
Set $\varepsilon_{j+1}=\delta(\varepsilon_j)$, $n_j=2^jn_0$ and
$\beta_j=2^j\beta_0$. Then for every even $L\ge n_0$ the ground state
is unique and
\begin{equation}
 \Delta_L\ \ge\ c:=\inf_{j\ge0}\frac1{\beta_j}
 \log\frac{1-3\varepsilon_j}{3\varepsilon_j}\ >\ 0 .
 \label{eq:c}
\end{equation}
\end{theorem}
\begin{proof}
By Theorem~\ref{thm:space}, the four moves of the doubling estimate in
Ref.~\cite{main} apply without change. Move~1 gives
$\Ps(n,2\beta)\ge(1-\varepsilon)^3$. Move~2 squares in space to get
$\Ps(2n,2\beta)\ge1-\delta(\varepsilon)$. Move~3 gives
$\Pt(4n,\beta)\ge(1-\varepsilon)^2(1-\delta)\ge(1-\varepsilon)^3$, which
requires $\delta\le\varepsilon$, i.e.\ $\varepsilon\le\varepsilon_*$.
Move~4 squares in time to get $\Pt(4n,2\beta)\ge1-\delta(\varepsilon)$.
Induction on $j$ gives
$\Ps(n_j,\beta_j),\Pt(2n_j,\beta_j)\ge1-\varepsilon_j$, with
$\varepsilon_j$ decreasing because $\delta(\varepsilon)<\varepsilon$ on
$(0,\varepsilon_*)$.
Since $3\varepsilon_j<1/2$, the interpolation and gap-recovery steps of
Ref.~\cite{main} give uniqueness and $e^{-\beta_j\Delta_L}\le
3\varepsilon_j/(1-3\varepsilon_j)$ whenever $n_j\le L\le2n_j$.
Every term of the infimum is positive; since
$\delta(\varepsilon)\le K\varepsilon^2$ near zero, $\varepsilon_j$ falls
doubly exponentially, so the terms converge to a positive limit.
\end{proof}
With $(n_0,\beta_0,\varepsilon_0)=(2304,784,1/100)$ this is the theorem of
Ref.~\cite{main}, with $c\ge(\log20)/784$.

\emph{Asymmetric defects and four-torus seed lifting.}
In Theorem~\ref{thm:criterion}, the initial spatial defect
$\varepsilon_{\mathrm s}\ge1-\Ps(n_0,\beta_0)$ and thermal defect
$\varepsilon_\tau\ge1-\Pt(2n_0,\beta_0)$ need not be equal. Tracking
$(\varepsilon_{\mathrm s,j},\varepsilon_{\tau,j})$ separately through Moves~1--4
gives the asymmetric recurrence
\begin{equation}
 \begin{aligned}
 \varepsilon_{\mathrm s,j+1}&=f\bigl(1-(1-\varepsilon_{\mathrm s,j})^2(1-\varepsilon_{\tau,j})\bigr),\\
 \varepsilon_{\tau,j+1}&=f\bigl(1-(1-\varepsilon_{\tau,j})^2(1-\varepsilon_{\mathrm s,j+1})\bigr),
 \end{aligned}
 \label{eq:asym}
\end{equation}
where $f(z)=z^2/[2(1-z)^2]$. Because the second line uses the updated spatial
defect $\varepsilon_{\mathrm s,j+1}$ from Move~2 rather than
$\varepsilon_{\mathrm s,j}$, the basin of attraction in the
$(\varepsilon_{\mathrm s},\varepsilon_\tau)$ plane extends well outside the
symmetric square $[0,\varepsilon_*)^2$, reaching $\varepsilon_{\mathrm s}\simeq0.2176$
as $\varepsilon_\tau\to0$ and $\varepsilon_\tau\simeq0.2006$ as
$\varepsilon_{\mathrm s}\to0$:
\begin{proofdiagram}
\begin{tikzpicture}[x=10.4cm,y=10.4cm,font=\scriptsize]
 % Asymmetric admissible region in (eps_s, eps_t) plane
 \fill[mpoblue!14]
   (0,0)--(0,.2006)--(.02,.1940)--(.04,.1860)--(.06,.1762)--(.08,.1645)
   --(.10,.1506)--(.12,.1341)--(.1273,.1273)--(.14,.1145)--(.16,.0915)
   --(.18,.0644)--(.20,.0326)--(.2176,0)--cycle;
 % Symmetric square [0, eps_*]^2
 \draw[black!45,densely dashed,line width=.4pt]
   (0,.1273)--(.1273,.1273)--(.1273,0);
 % Asymmetric boundary curve
 \draw[mpored,line width=.75pt]
   (0,.2006)--(.02,.1940)--(.04,.1860)--(.06,.1762)--(.08,.1645)
   --(.10,.1506)--(.12,.1341)--(.1273,.1273)--(.14,.1145)--(.16,.0915)
   --(.18,.0644)--(.20,.0326)--(.2176,0);
 % Axes
 \draw[black!55,line width=.4pt] (0,0)--(.235,0) node[right,black!75] {$\varepsilon_{\mathrm s}$};
 \draw[black!55,line width=.4pt] (0,0)--(0,.216) node[above,black!75] {$\varepsilon_\tau$};
 % Diagonal fixed point and certified asymmetric seeds
 \fill[black!70] (.1273,.1273) circle[radius=.0036]
   node[above right,black!75] {$\varepsilon_*$};
 \fill[mpoblue] (.0479,.1815) circle[radius=.004]
   node[above right,mpoblue!85!black] {spin-1 ($n_0{=}36$)};
 \fill[mpored] (.1856,.0511) circle[radius=.004]
   node[above right,mpored!85!black] {ladder, large-$D$, Ising};
 \node[black!55] at (.064,.055) {$[0,\varepsilon_*)^2$};
\end{tikzpicture}
\end{proofdiagram}
\noindent Certified seeds in both wings outside $[0,\varepsilon_*)^2$ close
under Eq.~\eqref{eq:asym}: the spin-1 seed $(0.0479,0.1815)$ of
Ref.~\cite{cert} lies in the upper wing, while the $J_\perp=5/2$ ladder
$(0.1856,0.0511)$, $D=3.206$ large-$D$ $(0.1916,0.0465)$, and $(g,h)=(1,1/20)$
Ising $(0.1840,0.0460)$ seeds lie in the right-hand wing.

Furthermore, reading the four-torus coupling identity of Ref.~\cite{main} to
solve for $\Pt(2n_0,\beta_0)$ yields
\begin{equation}
 \Pt(2n_0,\beta_0)=\frac{\Zt42}{\Ztsq41}
 =\frac{\Ps(n_0,2\beta_0)\,\Pt(n_0,\beta_0)^2}{\Ps(n_0,\beta_0)^2},
 \label{eq:seedlift}
\end{equation}
bounding $\Pt(2n_0,\beta_0)$ from length-$n_0$ and length-$2n_0$ partition
functions at $\beta_0$ and $2\beta_0$ without a ring of length $4n_0$. At $D=4$
(\path{models/spin1_largeD}), combining Eq.~\eqref{eq:seedlift} at $2n=12$ with
transfer squaring lifts the seed from $n_0=6$ to $n_0=12$ using only $6$- and
$12$-site spectra, raising the gap bound from $0.695$ to $0.899$.

\begin{corollary}[What (H3) excludes]\label{cor:obstruct}
Under (H1) and (H2), Eq.~\eqref{eq:H3} cannot hold at any scale
\begin{itemize}
\item[(i)] if the ground state is $k$-fold degenerate at the length $2n_0$
with $k\ge2$, because $\Pt\le w_0\le1/k$ there, a spectrum like \specbars{.49,.49,.02};
\item[(ii)] if $\liminf_L\Delta_L=0$ along even $L$, because
Theorem~\ref{thm:criterion} would contradict it.
\end{itemize}
\end{corollary}
Corollary~\ref{cor:obstruct} rules out gapless, symmetry-broken, and
asymptotically degenerate chains on the ring.

The converse is expected but not proved. In a gapped chain with a unique
ground state and correlation length $\xi$, one expects
$1-\Ps\sim e^{-n/\xi}$ and $1-\Pt\sim n\,e^{-\beta\Delta}$. Both vanish
along the doubling sequence \torusseq, so Eq.~\eqref{eq:H3} should hold at
some finite scale.

For a conformal critical chain the picture is different. The extensive free
energy cancels in each torus ratio, so $\Ps$ and $\Pt$ approach ratios of
the universal torus partition function at a fixed aspect ratio. These are
numbers below one, so the defect stalls instead of collapsing.

\section{Classification of Reflection-Positive Hamiltonians}\label{app:catalogue}
\emph{Bilinear spin couplings.} Consider spin~$S$ with a nearest-neighbor
coupling matrix $J$ (real symmetric) and a uniform field $\bm h$,
\begin{equation}
 H=\sum_j\bm S_j^{\mathsf T}J\,\bm S_{j+1}+\bm h\cdot\bm S_j .
 \label{eq:bilinear}
\end{equation}
\begin{proposition}\label{prop:xyz}
Suppose there is a symmetric orthogonal $O$ with $\det O=-1$ such that
$-JO\succeq0$ and $O\bm h=\bm h$. Then Eq.~\eqref{eq:bilinear} has the form
of Eq.~\eqref{eq:rpform} for every $S$. For diagonal
$J=\operatorname{diag}(J_x,J_y,J_z)$ with nonzero entries and $\bm h=0$, such
an $O$ exists if and only if an even number of $J_x,J_y,J_z$ are negative.
\end{proposition}
\begin{proof}
$O=-R$ with $R$ the identity or a $\pi$ rotation. Let $T$ be time reversal,
so $TST^{-1}=-S$, and put $\Theta=U_RT$. Then $\Theta\bm S\Theta^{-1}=O\bm S$.
Since $U_R$ commutes with $T$, $\Theta^2=U_R^2T^2=\pm1$. Factor
$-JO=\sum_kv_kv_k^{\mathsf T}$ with real $v_k$ and set
$C_k=v_k\cdot\bm S$. Then
$\sum_kC_k\otimes\Theta(C_k)=\sum\bm S^{\mathsf T}(-JO)O\bm S
=-\bm S^{\mathsf T}J\bm S$, and $\Theta(\bm h\cdot\bm S)=(O\bm h)\cdot\bm S$.
For diagonal $J$, $O=\operatorname{sign}(-J)$, so $\det O=-1$ counts the
ferromagnetic entries modulo two.
\end{proof}
\noindent The proposition sorts the standard spin chains as follows.
\begin{itemize}
\item \emph{Reflection positive:}
 \begin{itemize}
 \item the Heisenberg antiferromagnet for any $S$ ($O=-1$);
 \item XXZ chains $S^xS^x+S^yS^y+\Delta S^zS^z$ with $\Delta\ge0$;
 \item the XY model with any signs ($J_z=0$ frees $O_{zz}$);
 \item the transverse-field Ising chain with either sign ($O$ fixes the field
 axis), including a longitudinal field $-h\sum_j Z_j$ ($C_1=Z$, $o=-gX-hZ$).
 \end{itemize}
\item \emph{Not of this form:}
 \begin{itemize}
 \item the Heisenberg ferromagnet (three negative entries);
 \item XXZ chains with $\Delta<0$;
 \item the antiferromagnet in a uniform field, since $O=-1$ reverses $\bm h$.
 \end{itemize}
\end{itemize}

Sublattice rotations by $\pi$ change two signs at once and never alter the bond
parity, but they reverse a \emph{staggered} field $-h\sum_j(-1)^jS_j^z$,
\fieldstag: applying $U=2iS^y$ on odd sites maps the antiferromagnetic bond to
$-h_{\mathrm{rot}}=X\otimes X+G_y\otimes G_y+Z\otimes Z$ ($G_y=iS^y$) while
negating $S^z$ on odd sites, converting the staggered field into a uniform
reflection-even term $+h\sum_jS_j^z$ with a single-spin cell
(\path{models/heisenberg_staggered_field}). Similarly, in Ref.~\cite{main}
(Appendix~A) the open chain is built with opposite endpoint fields \fieldopp{}
because $\theta(-hS^z_1)=+hS^z_n$ is mirror symmetric, and the equal-field chain
\fieldsame{} is reached via $U_P$.

\begin{table*}[t]
\caption{Computer-certified uniform spectral gaps $\Delta_L=E_1(L)-E_0(L)$ established via the spacetime bootstrap (\texttt{models/}; $L,n_0$ count unit cells), proving ground-state uniqueness and $\Delta_L\ge\Delta_{\mathrm{cert}}$ for even $L$ in the stated range (odd $n$ in row~2). Lean~4 (\texttt{lean/CertifiedGaps/}, checked by \texttt{lean/scripts/check\_seeds.jl}) covers rows 1, 3, 4, 8, 9 (with $(1/4,1/20)$ proved unconditionally in \texttt{lean/SeedKernel/}), $\delta=9/10$ in row~5, $D\in\{4,1603/500\}$ at $n_0=6$ in row~6, and $J_\perp=4$ in row~7.\label{tab:certified}}
\begin{ruledtabular}
\begin{tabular}{llllll}
Model & Cell / RP form & Certified parameters & Seed $(n_0,\beta_0)$ & $\Delta_{\mathrm{cert}}$ & Range \\
\hline
Spin-1 Heisenberg & 1 site, $G_\alpha$ & isotropic ($J=1$) & $(36,49/4)\to(2304,784)$ & $(\log 20)/784$ & $L\ge2304$ \\
Spin-1 open~\cite{main} & 1 site, caps $v_\beta$ & fields $h=3/5$ & $(120,49)\to(1920,392)$ & $(\log 10)/392$ & odd $n\ge1921$ \\
Staggered Heisenberg & 1 site, $2iS^y$ & $h\in\{1,\tfrac34,\tfrac12,\tfrac7{20},\tfrac14\}$ & $(8,\,21/4\dots25/4)$ & $1.148\dots0.214$ & $L\ge8$ \\
Dimerized spin-$\tfrac12$ & dimer, swap $F$ & $\delta\in\{\tfrac12,\tfrac3{10},\tfrac14\}$ & $(4,\,6\dots29/4)$ & $0.534,\,0.173,\,0.058$ & $L\ge4$ cells \\
Dimerized spin-1 & dimer, $KK^*$ & $\delta\in\{\tfrac9{10},\tfrac7{10}\}$ & $(2,7/2),\,(3,17/4)^{\mathrm a}$ & $0.528,\,0.218$ & $L\ge2,\,4$ cells \\
Spin-1 large-$D$ & 1 site, $D(S^z)^2$ & $D=4;\,[3.2059998,3.206]$ & $(12,29/10)^{\mathrm b},(6,3.048)$ & $0.899;\,0.0010$ & $L\ge6$ \\
Heisenberg ladder & rung, $u=2iS^y$ & $J_\perp\in\{4,\tfrac72,3,\tfrac52,2.488\}$ & $(4,\,2\dots3)$ & $1.102\dots0.0015$ & $L\ge4$ rungs \\
Ising ladder & rung, $Z_1,Z_2$ & $g=3$ & $(4,3/2)$ & $0.636$ & $L\ge4$ rungs \\
Longitudinal Ising & 1 site, $-gX-hZ$ & 30 pts in $[\tfrac14,2]{\times}[\tfrac1{20},1]$ & $(6,\,2\dots97/4)$ & $0.0763\dots3.705$ & $L\ge6$ \\
\end{tabular}
\end{ruledtabular}
\footnotetext[1]{Odd cell seed $n_0=3$ ($6$ and $12$ spins) enabled by the positive cell-transfer factorization~\eqref{eq:cellpos} ($\delta=9/10$ uses the even seed $n_0=2$).}
\footnotetext[2]{Lifted from $n_0=6$ to $n_0=12$ via spatial transfer squaring and Eq.~\eqref{eq:seedlift}, with verified dense bridges at $L=6,8,10$.}
\end{table*}

\emph{Spin-1 bilinear--biquadratic chain.}
Consider $H=\sum_j\cos\vartheta\,\bm S_j\cdot\bm S_{j+1}
+\sin\vartheta\,(\bm S_j\cdot\bm S_{j+1})^2$. In the Cartesian basis the
quadrupoles $Q_{\alpha\beta}=S^\alpha S^\beta$ are real. Their antisymmetric
part is $\frac12\epsilon_{\alpha\beta\gamma}G_\gamma$, so
\begin{equation}
 -h=\bigl(\cos\vartheta-\tfrac12\sin\vartheta\bigr)\sum_\alpha G_\alpha\otimes\overline{G_\alpha}
 -\sin\vartheta\sum_{\alpha,\beta}Q^{+}_{\alpha\beta}\otimes\overline{Q^{+}_{\alpha\beta}},
 \label{eq:blbq}
\end{equation}
where $Q^+_{\alpha\beta}=\frac12(Q_{\alpha\beta}+Q_{\beta\alpha})$; the cross terms cancel between symmetric and antisymmetric parts.
\begin{proposition}\label{prop:blbq}
With $\Theta=K$ in the Cartesian basis, the bilinear--biquadratic chain is
in reflection-positive form if and only if
$-\pi+\arctan2\le\vartheta\le0$.
\end{proposition}
\begin{proof}
The condition holds if both coefficients in Eq.~\eqref{eq:blbq} are
nonnegative, which is the stated window. Conversely, the form exists exactly
when the realigned bond matrix is positive semidefinite after an on-site
shift; maximizing its least eigenvalue over all shifts gives a negative value
outside the window ($-0.087$ at $\vartheta=-0.66\pi$ and $-0.156$ at
$\vartheta=0.05\pi$; Sec.~\ref{app:checks}).
\end{proof}

\begin{proofdiagram}
\begin{tikzpicture}[x=1cm,y=1cm,font=\footnotesize]
 \def\R{1.15}
 % reflection-positive window [-pi+arctan 2, 0], in degrees [-116.57, 0]
 \fill[mpored!14] (0,0)--(0:\R) arc[start angle=0,end angle=-116.57,radius=\R]--cycle;
 \draw[black!35,line width=.4pt] (0,0) circle[radius=\R];
 \draw[mpored,line width=2.2pt] (0:\R) arc[start angle=0,end angle=-116.57,radius=\R];
 % method applies: Haldane part of the window, (-45, 0]
 \draw[mpoblue,line width=2.2pt] (0:{\R+.12}) arc[start angle=0,end angle=-45,radius={\R+.12}];
 % phase boundaries
 \foreach \a in {45,90,-45,-135}{\draw[black!55,line width=.4pt] (\a:{\R-.12})--(\a:{\R+.12});}
 \node[black!70] at (-22:{\R+.72}) {Haldane};
 \node[black!70] at (67.5:{\R+.38}) {critical};
 \node[black!70] at (180:{\R+.42}) {FM};
 \node[black!70] at (-90:{\R+.36}) {dimerized};
 % special points
 \foreach \a in {0,18.43,45,-45,-116.57}{\fill (\a:\R) circle[radius=1.4pt];}
 \node[anchor=west,font=\scriptsize] at (18.43:{\R+.08}) {AKLT};
 \node[anchor=west,font=\scriptsize] at ([yshift=-1pt]0:{\R+.2}) {Heisenberg};
 \node[anchor=south west,font=\scriptsize] at (45:{\R+.02}) {ULS};
 \node[anchor=north west,font=\scriptsize] at (-45:{\R+.14}) {TB};
 \node[anchor=north east,font=\scriptsize] at (-116.57:{\R+.02}) {$-\pi{+}\arctan2$};
 \draw[flow,black!40] (0,0)--(35:.55) node[pos=1.15,black!60,font=\scriptsize] {$\vartheta$};
\end{tikzpicture}
\end{proofdiagram}
The red arc is the reflection-positive window; the blue arc,
$-\pi/4<\vartheta\le0$, is where it meets the Haldane phase so
Theorem~\ref{thm:criterion} applies. Elsewhere in the window
Corollary~\ref{cor:obstruct} excludes it (degenerate dimerized ground states and
critical point TB). The Heisenberg chain sits on the edge of the window, while
AKLT and ULS lie outside it because the positive biquadratic coefficient has the
wrong sign, and optimizing over all on-site $\Theta$ finds none (Sec.~\ref{app:checks}).

\emph{Anisotropy, dimerized cells, and ladders.}
Any real symmetric on-site or intra-cell operator $o_j=\overline{o_j}$
preserves bond reflection positivity without requiring integrability:
\begin{itemize}
\item \emph{Single-ion anisotropy.} In the Cartesian spin-1 basis,
$o=D(S^z)^2=D\operatorname{diag}(1,1,0)$ is real diagonal and leaves
$-\bm S_j\cdot\bm S_{j+1}$ unchanged (\path{models/spin1_largeD}).
\item \emph{Dimerized chains.} Bond-alternating spin-$\tfrac12$ and spin-1
chains with couplings $1\pm\delta$ group into two-spin cells with swap $F=P$
under Eq.~\eqref{eq:rpcell} (\path{models/dimerized_heisenberg},
\path{models/spin1_dimerized}).
\item \emph{Two-leg ladders.} Grouping each two-spin rung into a four-state
site ($F=I$) places the rung coupling in $o_j$ and the leg couplings in
$-h_{\mathrm{cell}}=\sum_\alpha C_\alpha\otimes C_\alpha$, covering both the
Ising ladder ($C_{1,2}=Z\otimes I,\,I\otimes Z$; \path{models/two_leg_tfim})
and, under a checkerboard $\pi$-rotation $u=2iS^y$ on one spin per rung, the
spin-$\tfrac12$ Heisenberg ladder (\path{models/heisenberg_ladder}).
\end{itemize}
Reflection-positive fermion models~\cite{jaffe} and two-dimensional lattices
with reflection planes~\cite{fils,biskup} also satisfy (H2), with $X_\beta$
acting on a transverse column in 2D.

\section{Numerical Checks and Certified Models}\label{app:checks}
\emph{Reflection-positivity screening.}
Two floating-point tests screen Hamiltonians before algebraic certification:
\begin{itemize}
\item \emph{Gram test.} Tests Eq.~\eqref{eq:rp} on open chains of $4$ and $6$
sites, $\beta\in[0.3,3]$, reporting the least eigenvalue of
$A\mapsto\Tr[A\theta(A)e^{-\beta H}]$ relative to the largest, minimized over
$\Theta=VK$.
\item \emph{Realignment test} (\path{CertifiedGaps.rp_test}).
Checks Eq.~\eqref{eq:rpform} via the bond realignment matrix
$R[-h]_{(ik),(jl)}=-h_{(ij),(kl)}$ with an optimized on-site shift.
\end{itemize}
The screening results match Propositions~\ref{prop:xyz} and~\ref{prop:blbq}:
\begin{itemize}
\item \emph{Spin-$\tfrac12$ XYZ chain.} All eight sign patterns of
$(J_x,J_y,J_z)=(\pm1,\pm.7,\pm.4)$ follow the parity rule: zero or two
ferromagnetic entries give $\ge-10^{-12}$ after optimizing $V$, while the other
four give $-0.048$; the uniform-field antiferromagnet gives $-0.095$ and the
transverse-field Ising chain gives $\ge-10^{-12}$.
\item \emph{Spin-1, Cartesian $\Theta=K$.} The Heisenberg antiferromagnet and
$-(\bm S\cdot\bm S)^2$ pass; the ferromagnet ($-0.39$) and AKLT ($-0.04$) fail.
Over $\vartheta$, the test is nonnegative on $\vartheta/\pi\in[-.62,0]$
(Proposition~\ref{prop:blbq}) and negative on
$\vartheta/\pi\in\{-1,-.88,-.75\}\cup[.12,1]$.
\item \emph{Spin-1, every on-site $\Theta=VK$.} Minimizing over $V\in U(3)$
with twelve restarts leaves AKLT at $-0.040$, ULS at $-0.18$, and
$\vartheta=-0.75\pi$ and the ferromagnet at $-0.15$.
\end{itemize}
A failure excludes that on-site $\Theta$, not every block or cut.

\emph{Certified model suite and full-spectrum enclosures.}
Beyond the spin-1 Heisenberg chain (Ref.~\cite{cert}), the repository certifies
seven additional model families (Table~\ref{tab:certified}). For seed rings up
to $16$ spin-$\tfrac12$ or $12$ spin-1 sites, \path{CertifiedGaps.VerifiedSpectra}
assembles symmetry blocks (whose exact multiplicities sum to $d^n$) in interval
arithmetic, diagonalizes them in IEEE binary64, and encloses every eigenvalue
via polar-factor orthogonalization, Weyl's inequality, and a priori rounding
bounds ($\gamma_k=ku/(1-ku)$ with gradual underflow) before summing $\log Z$
and iterating Eq.~\eqref{eq:asym}. Centered operator-norm bounds on $\sum_j o_j$
extend point certificates to coupling intervals ($D\in[3.2059998,3.206]$ in
\path{models/spin1_largeD} and $|J_\perp-J_*|\le10^{-4}$ in
\path{models/heisenberg_ladder}, where independent-rung Weyl bounds fail for
all $n\ge2$, $\beta>0$, $0<J_\perp\le4$). Conversely, bounding the monotone
factors $W_L(\beta)=e^{\beta E_0(L)}Z_L(\beta)$ over interval coverings of
$(0,\infty)$ rules out all $\beta_0>0$ at $n_0\in\{2,4\}$ for the dimerized
spin-$\tfrac12$ chain at $\delta\le1/5$ and the Heisenberg ladder at
$J_\perp\in\{3/2,2\}$. At $(g,h)=(1/4,1/20)$ in the longitudinal-field Ising
chain, \path{lean/SeedKernel/} verifies the integer moments and seed bounds
in the Lean~4 kernel (\texttt{decide +kernel}), proving an unconditional gap.

\begin{thebibliography}{14}
\bibitem{main}
Companion note, \emph{A diagrammatic proof of the Haldane gap via spacetime
tensor networks} (\texttt{thermal-mpo-gap-note.tex}, 2026).
\bibitem{cert}
Companion note, \emph{How the initial estimate is certified}
(\texttt{thermal-mpo-gap-certificate.tex}, 2026).
\bibitem{os}
K. Osterwalder and R. Schrader,
\href{https://doi.org/10.1007/BF01645738}{Axioms for Euclidean Green's
functions}, Commun. Math. Phys. \textbf{31}, 83 (1973).
\bibitem{luscher}
M. L\"uscher, Construction of a selfadjoint, strictly positive transfer
matrix for Euclidean lattice gauge theories, Commun. Math. Phys.
\textbf{54}, 283 (1977).
\bibitem{fss}
J. Fr\"ohlich, B. Simon, and T. Spencer,
\href{https://doi.org/10.1007/BF01608557}{Infrared bounds, phase
transitions and continuous symmetry breaking}, Commun. Math. Phys.
\textbf{50}, 79 (1976).
\bibitem{jaffejanssens}
A. Jaffe and B. Janssens,
\href{https://doi.org/10.1007/s00220-015-2545-z}{Characterization of
reflection positivity: Majoranas and spins}, Commun. Math. Phys.
\textbf{346}, 1021 (2016).
\bibitem{dls}
F. J. Dyson, E. H. Lieb, and B. Simon,
\href{https://doi.org/10.1007/BF01106729}{Phase transitions in quantum spin
systems with isotropic and nonisotropic interactions},
J. Stat. Phys. \textbf{18}, 335 (1978).
\bibitem{fils}
J. Fr\"ohlich, R. Israel, E. H. Lieb, and B. Simon,
\href{https://doi.org/10.1007/BF01940327}{Phase transitions and reflection
positivity. I.}, Commun. Math. Phys. \textbf{62}, 1 (1978).
\bibitem{biskup}
M. Biskup, \emph{Reflection positivity and phase transitions in lattice spin
models}, in \emph{Methods of Contemporary Mathematical Statistical
Physics}, Lecture Notes in Math. \textbf{1970} (Springer, 2009).
\bibitem{speer}
E. R. Speer, Failure of reflection positivity in the quantum Heisenberg
ferromagnet, Lett. Math. Phys. \textbf{10}, 41 (1985).
\bibitem{jaffe}
A. Jaffe and F. L. Pedrocchi, Reflection positivity for Majoranas,
Ann. Henri Poincar\'e \textbf{16}, 189 (2015).
\bibitem{wuzhang}
C. Wu and S.-C. Zhang, Sufficient condition for absence of the sign problem
in the fermionic quantum Monte Carlo algorithm, Phys. Rev. B \textbf{71},
155115 (2005).
\bibitem{wei}
Z.-C. Wei, C. Wu, Y. Li, S. Zhang, and T. Xiang, Majorana positivity and the
fermion sign problem of quantum Monte Carlo simulations,
Phys. Rev. Lett. \textbf{116}, 250601 (2016).
\bibitem{li}
Z.-X. Li, Y.-F. Jiang, and H. Yao, Majorana-time-reversal symmetries: a
fundamental principle for sign-problem-free quantum Monte Carlo
simulations, Phys. Rev. Lett. \textbf{117}, 267002 (2016).
\end{thebibliography}
\end{document}
